\chapter{Cram\'er-Rao Lower Bound}
\label{cha:crlb}

This chapter presents a lower bound on the variance of an estimator, which offers a different perspective on optimality in the context of unbiased estimation. While previous chapters focused on characterizing optimal estimators—such as those found through complete statistics or conditioning—this ``lower bound" approach provides a complementary perspective by identifying fundamental barriers to estimation accuracy. 

The core idea is to establish a threshold that any estimator (under specific assumptions like unbiasedness) must respect; if an estimator's variance matches this lower bound, it is proven to be optimal as no better performance is possible. This bound also illustrates how estimation accuracy is impacted by the presence of nuisance parameters (i.e., parameters that are unknown but not of direct interest). For further reading, consider, e.g., Wellner (2018, chap.~3).

\begin{remark}
    The study of fundamental barriers is common in statistical literature. For instance, in minimax estimation, one seeks to prove lower bounds on the minimax risk. Proving a lower bound and subsequently exhibiting an estimator that achieves it (a matching upper bound) is a standard method to establish optimality. In this course, we focus on the simplest instance: the Cram\'er-Rao lower bound for unbiased estimation.
\end{remark}

Our starting point is the fundamental question:

\textbf{How well can one estimate a parameter?}

\paragraph{Setup}%
\label{par:Setup}
\begin{itemize}
    \item 
        \(X \sim P_{\bs{\theta}}\) where \(\bs{\theta} \in \Theta \subseteq \mathbb{R}^k\). We assume \(\Theta\) is an \textbf{open subset} to allow for differentiation with respect to the parameters.
    \item 
        Densities: 
        \(p_{\bs{\theta}}(x) = \frac{dP_{\bs{\theta}}}{d \nu}(x), x \in \mathcal{X}.\) 
        We assume the model is smoothly parametrized so that these densities are differentiable.
    \item 
        Target Parameter: \(\gamma(\bs{\theta}) \in \mathbb{R}^r\) for a differentiable map \(\gamma: \Theta \to \mathbb{R}^r\).
    \item 
        The Jacobian of \(\gamma\) is the \(r \times k\) matrix:
\end{itemize}

\[\dot{\gamma}(\bs{\theta}) = \left(\frac{\partial \gamma_i(\bs{\theta})}{\partial \theta_j}\right)_{ij}.\]

\begin{definition}[Bias]
    Suppose we have an estimator \(T: \mathcal{X} \to \mathbb{R}^r\) such that \(\mathsf{E}_{\bs{\theta}}[|T|] < \infty\) for all \(\bs{\theta} \in \Theta\). The \textbf{bias} of the estimator is defined as
    \[
        b(\bs{\theta}) = \mathsf{E}_{\bs{\theta}}[T] - \gamma(\bs{\theta}).
    \]
\end{definition}

\textbf{Question}: How small can \(\mathsf{Var}_{\bs{\theta}}[T]\) be in relationship to its bias? 

Specifically, for unbiased estimators, this question asks for the minimum achievable Mean Square Error (MSE).

\section{Score Function and Fisher
		Information}

\textbf{Definition 6.1.} Assuming existence of the derivatives and
expectations, we define

\begin{enumerate}[label=\roman*.]
	\item \textbf{Score function}

\[\dot{\ell}_{\bs{\theta}}(x) = \nabla_{\bs{\theta}} \log p_{\bs{\theta}}(x) = \begin{pmatrix} \frac{\partial}{\partial \bs{\theta}_{1}} \log p_{\bs{\theta}}(x) \\ \vdots \\ \frac{\partial}{\partial \bs{\theta}_{k}} \log p_{\bs{\theta}}(x) \end{pmatrix}.\]

    \item \textbf{Fisher Information} 

\[I(\theta) = \mathsf{Var}_{\theta}[\dot{\ell}_{\theta}(X)].\]
\end{enumerate}

In the sequel we derive a lower bound on \(\mathsf{Var}_{\theta}[T]\) , the
Cram\'er--Rao bound or also information inequality, in which the Fisher
information will play a key role.

\section{Cram\'er--Rao Bound}

If \(r = k = 1\), the Cram\'er--Rao (CR) bound will take the form

\[
	\mathsf{Var}_{\boldsymbol{\theta}}[T] \ge \frac{\left(\dot{\gamma}(\boldsymbol{\theta}) + \dot{b}(\boldsymbol{\theta})\right)^2}{\operatorname{I}(\boldsymbol{\theta})}, \quad \boldsymbol{\theta} \in \Theta.
\]

If \emph{T} is \textbf{unbiased}, then

\begin{equation}
	\label{eq:crb}
	\mathsf{Var}_{\boldsymbol{\theta}}[T] \ge \frac{\left(\dot{\gamma}(\boldsymbol{\theta})\right)^2}{\operatorname{I}(\boldsymbol{\theta})}, \quad \boldsymbol{\theta} \in \Theta.
\end{equation}

If \(r = 1\) and \(k>1\), the bound in \eqref{eq:crb} generalizes to

\[
	\mathsf{Var}_{\boldsymbol{\theta}}[T] \ge \dot{\gamma}(\boldsymbol{\theta}) \mathrm{I}^{-1}(\boldsymbol{\theta}) \dot{\gamma}(\boldsymbol{\theta})^{\top}.
\]

If \(r \ge 2\) , a lower bound on the \(r \times r\) covariance matrix
\(\mathsf{Var}_{\theta}[T]\) may be given in terms of the Löwner (or
positive semidefinite) ordering:

\[
	A \succeq B \iff A - B \succeq 0 \iff y^{\top}(A - B)y \geq 0 \quad \forall y \in \mathbb{R}^r.
\]

\paragraph{Assumptions:}%
\label{par:Assumptions}
\begin{enumerate}[label=(M\arabic*), ref=(M\arabic*)]
    \item\label{m1} \(\Theta\) is an open subset of \(\mathbb{R}^k\);
    \item \label{m2}\(\exists B \subset \mathcal{X}\) with \(\nu(B) = 0\) s.t.
	      \(\forall x \notin B : \theta \mapsto p_{\theta}(x)\) is differentiable;
      \item\label{m3} \(A := \{x : p_{\theta}(x) = 0\}\) does not depend on \(\theta\);\footnote{
        This is essentially the complement of the support
of \(\mathsf{P}_{\theta}\) , where the support is the smallest closed set \(C\) with
\(\mathsf{P}_{\theta}(C) = 1\) .}
\item\label{m4} \(I(\theta)\) is finite and positive definite for all
	      \(\theta \in \Theta\);
      \item \label{m5}The functions \(\bs{\theta} \mapsto \int p_{\theta}(x) dv(x)\) and
	      \(\bs{\theta} \mapsto \int T_i(x) p_{\theta}(x) dv(x)\) , \(1 \le i \le r\) ,
	      can be differentiated (wrt. \(\bs{\theta}\) ) under the integral sign. (Of
	      course, \(\int p_{\theta}(x) dv(x) \equiv 1\).)
\end{enumerate}

\begin{theo}[Cram\'er--Rao bound / Information inequality]
\label{theo:crb}
    
If (M1)-(M5) holds, then

\[
	\mathsf{Var}_{\boldsymbol{\theta}}[T(X)] \succeq \left(\dot{\gamma}(\boldsymbol{\theta}) + \dot{b}(\boldsymbol{\theta})\right) \mathbf{I}^{-1}(\boldsymbol{\theta}) \left(\dot{\gamma}(\boldsymbol{\theta}) + \dot{b}(\boldsymbol{\theta})\right)^{\top}, \quad \boldsymbol{\theta} \in \Theta.
\]

\end{theo}

\begin{proof}[]
    
\[
	\mathsf{Var}_{\boldsymbol{\theta}}\bigg[\begin{pmatrix} T(X) \\ \dot{\ell}_{\boldsymbol{\theta}}(X) \end{pmatrix}\bigg] = \begin{pmatrix} \mathsf{Var}_{\boldsymbol{\theta}}[T(X)] & \mathsf{Cov}_{\boldsymbol{\theta}}[T(X), \dot{\ell}_{\boldsymbol{\theta}}(X)] \\ \mathsf{Cov}_{\boldsymbol{\theta}}[T(X), \dot{\ell}_{\boldsymbol{\theta}}(X)] & \mathsf{Var}_{\boldsymbol{\theta}}[\dot{\ell}_{\boldsymbol{\theta}}(X)] \end{pmatrix}
\]

By definition, \(\mathsf{Var}_{\theta}[\dot{\ell}_{\theta}(X)] = I(\theta)\) .
Under our assumptions:

\begin{equation}
	\label{eq:fact}
	\boxed{\text{Fact: }
	\mathsf{E}_{\boldsymbol{\theta}}[\dot{\ell}_{\boldsymbol{\theta}}(X)] = 0}
\end{equation}

Indeed, for \(j = 1, \ldots, k\):

\[
    \begin{split} \mathsf{E}_{\bs{\theta}}[\dot{\ell}_{\bs{\theta},j}(X)] &= \int \left(\frac{\partial}{\partial \bs{\theta}_{j}} \log p_{\bs{\theta}}(x)\right) p_{\bs{\theta}}(x) \, \mathrm{d}\nu(x) \\ &= \int \frac{\partial}{\partial \bs{\theta}_{j}} p_{\bs{\theta}}(x) \, \mathrm{d}\nu(x) \stackrel{\ref{m5}}{=} \frac{\partial}{\partial \bs{\theta}_{j}} \int p_{\bs{\theta}}(x) \, \mathrm{d}\nu(x) \\ &= \frac{\partial}{\partial \bs{\theta}_{j}} 1 = 0. \end{split}
\]

It also holds that
\(\mathsf{Cov}_{\pmb{\theta}}[T(X), \dot{\ell}_{\pmb{\theta}}(X)] = \dot{\gamma}(\pmb{\theta}) + \dot{b}(\pmb{\theta})\)
as seen from

\[
	\begin{aligned} \mathsf{Cov}_{\bs{\theta}} \big[ T_{i}(X), \dot{\ell}_{\bs{\theta}, j}(X) \big] &\stackrel{\eqref{eq:fact} }{=} \mathsf{E}_{\bs{\theta}} \big[ T_{i}(X) \dot{\ell}_{\bs{\theta}, j}(X) \big] \\ &= \int T_{i}(x) \Big( \frac{\partial}{\partial \bs{\theta}_{j}} \log p_{\bs{\theta}}(x) \Big) p_{\bs{\theta}}(x) \, \mathrm{d}\nu(x) \\ &= \int \frac{\partial}{\partial \bs{\theta}_{j}} T_{i}(x) p_{\bs{\theta}}(x) \, \mathrm{d}\nu(x) \\ &\stackrel{\ref{m5}}{=} \frac{\partial}{\partial \bs{\theta}_{j}} \mathsf{E}_{\bs{\theta}} \big[ T_{i}(X) \big] = \big( \dot{\gamma}(\bs{\theta}) + \dot{b}(\bs{\theta}) \big)_{ij}. \end{aligned}
\]

Therefore,

\[\mathsf{Var}_{\boldsymbol{\theta}} \bigg[ \begin{pmatrix} T(X) \\ \dot{\ell}_{\boldsymbol{\theta}}(X) \end{pmatrix} \bigg] = \begin{pmatrix} \mathsf{Var}_{\boldsymbol{\theta}} [T(X)] & \dot{\gamma}(\boldsymbol{\theta}) + \dot{b}(\boldsymbol{\theta}) \\ \left( \dot{\gamma}(\boldsymbol{\theta}) + \dot{b}(\boldsymbol{\theta}) \right)^\top & \mathrm{I}(\boldsymbol{\theta}) \end{pmatrix}.\]

Since this covariance matrix is positive semidefinite and \(I(\theta)\)
is invertible, it follows that the Schur complement (see \autoref{def:schur})
of the Fisher information is equal to

\[
	\mathsf{Var}_{\boldsymbol{\theta}}[T(X)] - (\dot{\gamma}(\boldsymbol{\theta}) + \dot{b}(\boldsymbol{\theta}))I^{-1}(\boldsymbol{\theta})(\dot{\gamma}(\boldsymbol{\theta}) + \dot{b}(\boldsymbol{\theta}))^{\top} \succeq 0,
\]


and positive semidefinite, which completes the proof.
\end{proof}

\subsubsection*{Schur Complement}%
\label{ssub:schur}
Consider a partitioned matrix

\[
A =
\begin{pNiceMatrix}[first-row, first-col]
      & k & m \\
k & A_{11} & A_{12} \\
m & A_{21} & A_{22}
\end{pNiceMatrix}
\]

in which the \(m \times m\) block \(A_{22}\) is invertible.

\begin{definition}[Schur complement]
    \label{def:schur}
The Schur complement of \(A_{22}\) in A is

\[A_{11.2} := A_{11} - A_{12}A_{22}^{-1}A_{21}.\]
\end{definition}

\begin{lemma}
If \(A \succeq 0\) and in particular A symmetric,
then \(A_{11.2} \succeq 0\) . And if \(A \succ 0\) , then
\(A_{11.2} \succ 0\).\footnote{Notation: \(\succeq \) 0 means positive semidefinite,
and \(\succ\) 0 means positive definite.}
\end{lemma}

\begin{proof}[]
By symmetry of A,

\[\begin{pmatrix} \mathbf{I} & -A_{12}A_{22}^{-1} \\ \mathbf{0} & \mathbf{I} \end{pmatrix}^{\mathsf{T}} = \begin{pmatrix} \mathbf{I} & \mathbf{0} \\ -A_{22}^{-1}A_{21} & \mathbf{I} \end{pmatrix},\]

and, thus,

\begin{equation}
    \label{eq:mat}
	\begin{pmatrix} I & -A_{12}A_{22}^{-1} \\ 0 & I \end{pmatrix} \begin{pmatrix} A_{11} & A_{12} \\ A_{21} & A_{22} \end{pmatrix} \begin{pmatrix} I & -A_{12}A_{22}^{-1} \\ 0 & I \end{pmatrix}^{\top} = \begin{pmatrix} A_{11.2} & 0 \\ 0 & A_{22} \end{pmatrix}.
\end{equation}

For any \(\bs{y} \in \mathbb{R}^k\) , let
\(\bs{z} = \begin{pmatrix} \bs{I} & -A_{12}A_{22}^{-1} \\ \bs{0} & \bs{I} \end{pmatrix}^{\top} \begin{pmatrix} \bs{y} \\ \bs{0} \end{pmatrix} \in \mathbb{R}^{k+m}\)
. Then

\[
	\begin{split} \bs{0} &\leq \bs{z}^{\top} A \bs{z} \\ &= \begin{pmatrix} \bs{y}^{\top} & \bs{0} \end{pmatrix} \begin{pmatrix} \bs{I} & -A_{12} A_{22}^{-1} \\ \bs{0} & \bs{I} \end{pmatrix} \begin{pmatrix} A_{11} & A_{12} \\ A_{21} & A_{22} \end{pmatrix} \begin{pmatrix} \bs{I} & -A_{12} A_{22}^{-1} \\ \bs{0} & \bs{I} \end{pmatrix}^{\top} \begin{pmatrix} \bs{y} \\ \bs{0} \end{pmatrix} \\ &= \begin{pmatrix} \bs{y}^{\top} & \bs{0} \end{pmatrix} \begin{pmatrix} A_{11.2} & \bs{0} \\ \bs{0} & A_{22} \end{pmatrix} \begin{pmatrix} \bs{y} \\ \bs{0} \end{pmatrix} \\ &= \bs{y}^{\top} A_{11.2} \bs{y}, \end{split}
\]

for all \(y \in \mathbb{R}^k\) , and thus \(A_{11.2} \succeq 0\) .

If \(A \succ 0\), then \(\bs{z}^{\top} A \bs{z} > 0\) for
all \(\bs{z} \neq 0\) . Since \(\bs{z} \neq 0\) implies
\(\bs{y} \neq 0\) and hence
\(\bs{y}^{\top} A_{11.2} \bs{y} > 0\) for all
\(\bs{y} \neq 0\) , which shows \(A_{11.2} \succ 0\) .
\end{proof}


\subsubsection{Inverting a Partitioned Matrix}%
\label{ssub:ipm}

\begin{lemma}
Assuming all inverses exists:

\[A^{-1} = \begin{pmatrix} A_{11} & A_{12} \\ A_{21} & A_{22} \end{pmatrix}^{-1} = \begin{pmatrix} (A_{11.2})^{-1} & -A_{11.2}^{-1}A_{12}A_{22}^{-1} \\ -A_{22}^{-1}A_{21}A_{11.2}^{-1} & (A_{22.1})^{-1} \end{pmatrix},\]

where
\(A_{11,2}^{-1}A_{12}A_{22}^{-1} = A_{11}^{-1}A_{12}A_{22,1}^{-1}\) .
\end{lemma}

\begin{proof}[]
Check \(A \cdot A^{-1} = I\) or invert \eqref{eq:mat}  and rearrange
the pieces.
\end{proof}

\section{More on Fisher Information}%
\label{sec:mfi}

\subsection{Fisher Information and Curvature of the Log--likelihood Function}%
\label{sub:ficlogf}
Under \ref{m1}--\ref{m5},

\[\mathrm{I}(\bs{\theta}) := \mathsf{Var}_{\bs{\theta}} \big[ \dot{\ell}_{\bs{\theta}}(X) \big] = \mathsf{E}_{\bs{\theta}} \big[ \dot{\ell}_{\bs{\theta}}(X) \dot{\ell}_{\bs{\theta}}(X)^{\top} \big]\]

because \(\mathsf{E}_{\bs{\theta}}[\dot{\ell}_{\bs{\theta}}(X)] = 0\) .

If in addition it holds that \textbf{(M6)} 

\[
	\int p_{\theta}(x) d\nu(x)\quad \quad  \text{can be differentiated twice under the
integral sign,}
\]
then,

\[\mathrm{I}(\boldsymbol{\theta}) = -\mathsf{E}_{\boldsymbol{\theta}} \big[ \ddot{\boldsymbol{\ell}}_{\boldsymbol{\theta}}(\boldsymbol{X}) \big]\]

where
\(\ddot{\ell}_{\bs{\theta}}(x) = \left(\frac{\partial^2}{\partial \bs{\theta}_i \partial \bs{\theta}_j} \log p_{\bs{\theta}}(x)\right)_{ij}\)
is the Hessian of the log--likelihood function.

\textbf{NB:}  The Fisher information captures expected curvature of the
log--likelihood function!

\begin{proof}[]
    
We have that

\[\begin{split} \frac{\partial^2}{\partial \bs{\theta}_i \partial \bs{\theta}_j} \log p_{\bs{\theta}}(x) &= \frac{1}{p_{\bs{\theta}}(x)} \frac{\partial^2}{\partial \bs{\theta}_i \partial \bs{\theta}_j} p_{\bs{\theta}}(x) - \frac{1}{p_{\bs{\theta}}(x)^2} \left( \frac{\partial}{\partial \bs{\theta}_i} p_{\bs{\theta}}(x) \right) \left( \frac{\partial}{\partial \bs{\theta}_j} p_{\bs{\theta}}(x) \right) \\ &= \frac{1}{p_{\bs{\theta}}(x)} \frac{\partial^2}{\partial \bs{\theta}_i \partial \bs{\theta}_j} p_{\bs{\theta}}(x) - \left( \frac{\partial}{\partial \bs{\theta}_i} \log p_{\bs{\theta}}(x) \right) \left( \frac{\partial}{\partial \bs{\theta}_j} \log p_{\bs{\theta}}(x) \right). \end{split}\]

Take expectations,

\[\begin{split} \mathsf{E}_{\bs{\theta}} & \left[ \frac{\partial^2}{\partial \bs{\theta}_i \partial \bs{\theta}_j} \log p_{\bs{\theta}}(X) \right] \\ & = \int \frac{\partial^2}{\partial \bs{\theta}_i \partial \bs{\theta}_j} p_{\bs{\theta}}(x) \mathrm{d} \nu(x) - \mathsf{E}_{\bs{\theta}} \left[ \left( \frac{\partial}{\partial \bs{\theta}_i} \log p_{\bs{\theta}}(X) \right) \left( \frac{\partial}{\partial \bs{\theta}_j} \log p_{\bs{\theta}}(X) \right) \right] \\ & = 0 - \left( \mathsf{E}_{\bs{\theta}} \left[ \dot{\ell}_{\bs{\theta}}(X) \dot{\ell}_{\bs{\theta}}(X)^{\top} \right] \right)_{ij}. \end{split}\]
\end{proof}


\subsection{Fisher Information and Random Samples}%
\label{sub:firs}
\begin{pro}[]
    
Suppose \(\bs{X} = (X_1, ..., X_n)\) with
\(X_1, ..., X_n \overset{i.i.d.}{\sim} P_{\theta}, \theta \in \Theta\) .
Let \(I_n(\theta) = \mathsf{Var}_{\theta} [\dot{\ell}_{\theta}(X)]\) be
the Fisher information. Then under \ref{m1}--\ref{m5}:

\[
	I_n(\boldsymbol{\theta}) = nI_1(\boldsymbol{\theta}).
\]
\end{pro}

\begin{proof}[]
By the assumed independence, it holds that

\[\dot{\ell}_{\bs{\theta}}(\bs{X}) = \nabla_{\bs{\theta}} \sum_{i=1}^{n} \log p_{\bs{\theta}}(X_i) = \sum_{i=1}^{n} \dot{\ell}_{\bs{\theta}}(X_i).\]

Taking variances gives

\[\mathrm{I}_n(\bs{\theta}) = \mathsf{Var}_{\boldsymbol{\bs{\theta}}} \big[ \dot{\ell}_{\boldsymbol{\bs{\theta}}}(\bs{X}) \big] = \sum_{i=1}^n \mathsf{Var}_{\boldsymbol{\bs{\theta}}} \big[ \dot{\ell}_{\boldsymbol{\bs{\theta}}}(X_i) \big] = n \mathrm{I}_1(\boldsymbol{\bs{\theta}}) \,.\]
\end{proof}

\section{Examples}%
\label{sec:Examples}

\begin{ex}[Gaussian model]
    \label{ex:gm}
    
Let \(X_1, \ldots, X_n\) be i.i.d.
\(\mathcal{N}(\mu, \sigma^2), \theta = (\mu, \sigma^2) \in \mathbb{R} \times (0, \infty)\)
. Some calculus gives

\[I_1(\mu, \sigma^2) = \begin{pmatrix} \frac{1}{\sigma^2} & 0\\ 0 & \frac{1}{2\sigma^4} \end{pmatrix} \text{ and } I_1^{-1}(\mu, \sigma^2) = \begin{pmatrix} \sigma^2 & 0\\ 0 & 2\sigma^4 \end{pmatrix}.\]

\emph{CR} bound for unbiased estimation of \(\mu\) :

\[\mathsf{Var}_{\theta}[T] \ge \frac{\sigma^2}{n},\] 

achieved by \(T(\bs{X}) = \overline{X}_n\) .

CR bound for unbiased estimation of \(\sigma^2\):

\[\mathsf{Var}_{\boldsymbol{\theta}}[T] \ge \frac{2\sigma^4}{n},\]


not achieved because the UMVUE is the sample variance \(s^2\) with

\[\mathsf{Var}_{\boldsymbol{\theta}}[s^2] = \frac{2\sigma^4}{n-1}.\]

But note that bound is achieved asymptotically for \(n \to \infty\) .
\end{ex}

\begin{ex}[Reparametrization of Gaussian model]
    
In \autoref{ex:gm}, \(I_1(\mu, \sigma^2)\) is diagonal, which need not stay
true if we reparametrize.

Let \[\nu = \mathsf{E}_{\mu,\sigma^2}[X_1^2].\]  
Then

\[\begin{pmatrix} \mu \\ \sigma^2 \end{pmatrix} = \begin{pmatrix} \mu \\ \nu - \mu^2 \end{pmatrix}.\]

and {[}check!{]}

\[\begin{split} \mathrm{I}_{1}(\mu,\nu) &= \begin{pmatrix} \frac{\mu^{2}+\nu}{(\mu^{2}-\nu)^{2}} & \frac{-\mu}{(\mu^{2}-\nu)^{2}} \\ \frac{-\mu}{(\mu^{2}-\nu)^{2}} & \frac{1}{2(\mu^{2}-\nu)^{2}} \end{pmatrix}, \\ \mathrm{I}_{1}^{-1}(\mu,\nu) &= \begin{pmatrix} \nu-\mu^{2} & 2\mu(\nu-\mu^{2}) \\ 2\mu(\nu-\mu^{2}) & 2(\nu^{2}-\mu^{4}) \end{pmatrix}. \end{split}\]

Since \(\nu - \mu^2 = \sigma^2\) , the sample mean \(\overline{X}_n\)
still achieves the CR bound obtained from this parametrization (as it
should\ldots).

Moreover, \(\hat{\nu} = \frac{1}{n} \sum_{i=1}^{n} X_i^2\) , with
\(\mathsf{Var}_{(\mu,\nu)}[\hat{\nu}] = \frac{1}{n} \mathsf{Var}_{(\mu,\nu)}[X_1^2]\)
, can be shown to achieve the CR bound for unbiased estimation of
\(\nu\) .
\end{ex}

\begin{ex}[Exponential family in natural parametrization]
    
Consider an observation X that follows a distribution \(\mathsf{P}_{\eta}\) from
an exponential family in canonical form with density

\[p_{\eta}(x) = \exp \{ \langle \eta, T(x) \rangle - A(\eta) \} h(x), \quad x \in \mathcal{X}.\]

Then \(\dot{\ell}_{\eta}(x) = T(x) - \nabla_{\eta} A(\eta)\) and

\[
	\begin{split} \mathrm{I}(\eta) &= \mathsf{Var}_{\eta}[\dot{\ell_{\eta}}(X)] = \mathsf{Var}_{\eta}[T(X)] \\ &= D_2 A(\eta) = -\mathsf{E}_{\eta}[\ddot{\ell_{\eta}}(X)]. \end{split}
\]

Consider estimating the vector of mean parameters
\(\gamma(\eta) = \mathsf{E}_{\eta}[T(X)]\) . Then
\(\dot{\gamma}(\eta) = D_2 A(\eta)\) and the CR bound takes the form:

\begin{align*}
    \dot{\gamma}(\eta)\bs{I}^{-1}(\eta)\dot{\gamma}(\eta)^{\top} &= D_2 A(\eta) \left(D_2 A(\eta)\right)^{-1} D_2 A(\eta)\\
                                                                 &= D_2 A(\eta) 
                                                               \\&= \mathsf{Var}_{\eta}[T(X)].
\end{align*}

We observe that \emph{T} achieves CR bound for unbiased estimation of
\(\gamma(\eta) = \mathsf{E}_{\eta}[T(X)]\).

This provides a second proof that \emph{T} is UMVUE (besides using the
theorem of Lehmann-Scheffé).
\end{ex}

\begin{ex}[Gamma model]
    
Let \(X_1,\ldots,X_n\) be i.i.d. Gamma \((\alpha,\beta),\)
\(\bs{\theta}=(\alpha,\beta)\in(0,\infty)^2.\) The Lebesgue density is

\[p_{\bs{\theta}}(x) = \frac{\beta^{\alpha}}{\Gamma(\alpha)} x^{\alpha - 1} e^{-\beta x} \bs{1}_{(0,\infty)}(x)\]

\[= \exp\left\{ \left\langle \begin{pmatrix} \alpha - 1 \\ -\beta \end{pmatrix}, \begin{pmatrix} \log x \\ x \end{pmatrix} \right\rangle \right.\]

\[\left.- \left( \log \Gamma(\alpha) - \alpha \log \beta \right) \right\} \bs{1}_{(0,\infty)}(x).\]

The Fisher information can be shown to be

\[I_1(\alpha,\beta) = \begin{pmatrix} \psi^{(1)}(\alpha) & -\frac{1}{\beta} \\ -\frac{1}{\beta} & \frac{\alpha}{\beta^2} \end{pmatrix},\]

with inverse

\[I_1^{-1}(\alpha,\beta) = \frac{1}{\alpha\psi^{(1)}(\alpha) - 1} \begin{pmatrix} \alpha & \beta \\ \beta & \beta^2\psi^{(1)}(\alpha) \end{pmatrix},\]

where \(\psi^{(1)}(\alpha) = \frac{d^2}{d\alpha^2} \log \Gamma(\alpha)\)
.

Then

\[I_n(\alpha, \beta) = nI_1(\alpha, \beta) \quad \quad \text{and}\quad \quad 
I_n(\alpha, \beta)^{-1} = \frac{1}{n}I_1(\alpha, \beta)^{-1}\] .
\end{ex}

\begin{remark}[]
    \leavevmode
\begin{itemize}
	\item
	      Cram\'er--Rao bound need not be achieved by UMVUE.
	\item
	      Bound is attained for mean parameters of exponential families.
	\item
	      Under a differentiability assumption, achievement of Cram\'er--Rao bound
	      requires an exponential family (and consideration of their mean
	      parameters).
	\item
	      We will later see that MLE achieves Cram\'er-Rao bound asymptotically.
	\item
	      Differentiability requirements can be weakened using the concept of
	      differentiability in quadratic mean; see, e.g., van der Vaart (1998,
	      sec.~7.2) for further reading if you are interested.
\end{itemize}
\end{remark}

\section{Nuisance Parameters}
\label{sec:nuisance-parameters}

In many statistical problems, the parameter vector $\bs{\theta}$ contains components of primary interest, alongside other components that are unknown but of secondary interest. These secondary parameters are often required to fully specify the model but are not the direct target of the estimation. We refer to these as \textbf{nuisance parameters}.

\begin{ex}[Speed of Light]
    Consider measuring the speed of light (a physical constant) using a device with unknown accuracy. We might model the measurements using a normal distribution $N(\mu, \sigma^2)$. 
    \begin{itemize}
        \item The mean $\mu$ represents the speed of light. This is the \textbf{parameter of interest}.
        \item The variance $\sigma^2$ represents the measurement error spread. This is unknown, since we did not collect data specifically to learn about the device's accuracy. Thus, $\sigma^2$ is a \textbf{nuisance parameter}.
    \end{itemize}
\end{ex}

A fundamental question arises: 

\textit{What is the "price" we pay in terms of estimation accuracy for not knowing the nuisance parameter?} 

We can answer this using the Cram\'er-Rao bound.

\subsection{Partitioned Information}

Suppose the model is given by
\[
\mathcal{P} = \{ \mathsf{P}_{\bs{\theta}} : \bs{\theta} = (\bs{\theta}_1, \bs{\theta}_2) \in \Theta_1 \times \Theta_2 \},
\]
where $\bs{\theta}_1$ is the parameter of interest and $\bs{\theta}_2$ is the nuisance parameter. The Fisher information matrix can be naturally partitioned into four blocks:

\[
\bs{I}(\bs{\theta}) = \begin{pmatrix} \bs{I}_{11}(\bs{\theta}) & \bs{I}_{12}(\bs{\theta}) \\ \bs{I}_{21}(\bs{\theta}) & \bs{I}_{22}(\bs{\theta}) \end{pmatrix}.
\]

Here, $\bs{I}_{11}$ corresponds to partial derivatives with respect to $\bs{\theta}_1$, $\bs{I}_{22}$ to $\bs{\theta}_2$, and the off-diagonal blocks represent cross-covariances.

Consider the estimation of the target parameter \[
	\gamma(\bs{\theta}) = \bs{\theta}_1
    = \begin{pmatrix} I & 0 \end{pmatrix} \begin{pmatrix} \bs{\theta}_{1}\\\bs{\theta}_{2} \end{pmatrix}.
\]The Jacobian of this mapping is the projection matrix:
\[
\dot{\gamma}(\bs{\theta}) = \begin{pmatrix} \bs{I} & \bs{0} \end{pmatrix},
\]
where $\bs{I}$ is the identity matrix corresponding to the dimension of $\bs{\theta}_1$.

The Cram\'er-Rao (CR) bound at $\bs{\theta}_0$ for an unbiased estimator of $\bs{\theta}_1$ is given by the top-left block of the inverse information matrix:

\[
\begin{split} 
\text{CR Bound} &= \dot{\gamma}(\bs{\theta}_0) \bs{I}^{-1}(\bs{\theta}_0) \dot{\gamma}(\bs{\theta}_0)^\top \\
&= \begin{pmatrix} \bs{I} & \bs{0} \end{pmatrix} \bs{I}^{-1}(\bs{\theta}_0) \begin{pmatrix} \bs{I} \\ \bs{0} \end{pmatrix} \\ 
&= \left( \bs{I}^{-1}(\bs{\theta}_0) \right)_{11}.
\end{split}
\]

Using the formula for the inverse of a block matrix, this quantity is the inverse of the \textbf{Schur complement}:

\[
\left( \bs{I}^{-1}(\bs{\theta}_0) \right)_{11} = \left( \bs{I}_{11}(\bs{\theta}_0) - \bs{I}_{12}(\bs{\theta}_0) \bs{I}_{22}(\bs{\theta}_0)^{-1} \bs{I}_{21}(\bs{\theta}_0) \right)^{-1}.
\]

We define the \textbf{effective information} for $\bs{\theta}_1$ in the presence of unknown $\bs{\theta}_2$ as:
\[
\bs{I}_{11.2}(\bs{\theta}_0) := \bs{I}_{11}(\bs{\theta}_0) - \bs{I}_{12}(\bs{\theta}_0)\bs{I}_{22}(\bs{\theta}_0)^{-1}\bs{I}_{21}(\bs{\theta}_0).
\]
Thus, the CR bound is $\bs{I}_{11.2}(\bs{\theta}_0)^{-1}$.

\subsection{Comparison: Known vs. Unknown Nuisance Parameters}

To understand the cost of ignorance, consider a submodel where the nuisance parameter is known to be $\bs{\theta}_2 = \bs{\theta}_{20}$:
\[
\mathcal{P}(\bs{\theta}_{20}) = \{ \mathsf{P}_{\bs{\theta}} : \bs{\theta} = (\bs{\theta}_1, \bs{\theta}_{20}), \bs{\theta}_1 \in \Theta_1 \}.
\]
In this scenario, estimating $\bs{\theta}_1$ relies only on the partial derivatives with respect to $\bs{\theta}_1$. The Fisher information for this submodel is simply the top-left block $\bs{I}_{11}(\bs{\theta}_0)$. Consequently, the CR bound becomes $\bs{I}_{11}(\bs{\theta}_0)^{-1}$.

We can now compare the information available in both scenarios:
\begin{itemize}
    \item \textbf{With $\bs{\theta}_2$ known:} Information is $\bs{I}_{11}(\bs{\theta}_0)$.
    \item \textbf{With $\bs{\theta}_2$ unknown:} Information is $\bs{I}_{11.2}(\bs{\theta}_0) = \bs{I}_{11}(\bs{\theta}_0) - \bs{I}_{12}\bs{I}_{22}^{-1}\bs{I}_{21}$.

Since $\bs{I}(\bs{\theta})$ is a covariance matrix, it is symmetric and positive semi--definite. The term being subtracted, $\bs{I}_{12}\bs{I}_{22}^{-1}\bs{I}_{21}$, is of the form $A B A^\top$ with $B$ positive definite, meaning the subtraction term is positive semi--definite. Therefore:
\[
\bs{I}_{11}(\bs{\theta}_0) \succeq \bs{I}_{11.2}(\bs{\theta}_0).
\]
\begin{remark}[The Information Inequality]
    In general, not knowing the nuisance parameter $\bs{\theta}_2$ strictly reduces the information available for estimating $\bs{\theta}_1$ (and thus increases the variance bound), unless:
    \[
    \bs{I}_{12}(\bs{\theta}_0) = \bs{0}.
    \]
    If the Fisher information matrix is block diagonal, the parameters are orthogonal. in this ``lucky" case, there is no price to pay for not knowing the nuisance parameter. For example, when estimating the mean $\mu$ of a normal distribution, the sample mean $\ol{X}$ is the optimal estimator regardless of whether the variance $\sigma^2$ is known or unknown.
\end{remark}
\end{itemize}

\section{Reparametrization}
\label{sec:reparametrization}

Consider a statistical model \(\cc{P} = \{ \mathsf{P}_{\bs{\theta}} : \bs{\bs{\theta}} \in \bs{\theta} \}\) with Fisher information \(I(\bs{\theta}_0)\) at \(\bs{\theta}_0 \in \bs{\theta}\).

In practice, different software or researchers may parametrize the same distribution differently (e.g., using variance versus precision, or rate versus scale). Suppose we reparametrize the model using a \textbf{diffeomorphism}\footnote{A differentiable map with a differentiable inverse.} that maps each new parameter \(\bs{\lambda} \in \Lambda\) to a point \(\bs{\theta}(\bs{\lambda}) \in \bs{\theta}\).

\begin{pro}
    \label{pro:crb}
    The Fisher information of the reparametrized model \(\cc{P} = \{\mathsf{P}_{\bs{\theta}(\bs{\lm})} : \bs{\lm} \in \bs{\lm}\}\) at the point \(\bs{\lm}_0\) with \(\bs{\theta}(\bs{\lm}_0) = \bs{\theta}_0\) is equal to
    \[
    \bs{I}(\bs{\lm}_0) = D_{\bs{\lm}} \bs{\theta}(\bs{\lm})^{\top} \Big|_{\bs{\lm} = \bs{\lm}_0} \cdot \bs{I}(\bs{\theta}_0) \cdot D_{\bs{\lm}} \bs{\theta}(\bs{\lm}) \Big|_{\bs{\lm} = \bs{\lm}_0},
    \]
    where \(D_{\bs{\lm}}\bs{\theta}(\bs{\lm})\) denotes the Jacobian matrix of the map \(\bs{\lm} \mapsto \bs{\theta}(\bs{\lm})\).
\end{pro}

\begin{proof}
    By the chain rule, the reparametrized model has the score function:
    \[
    \dot{\ell}_{\bs{\lm}}(X) \equiv \nabla_{\bs{\lm}} \log p_{\bs{\theta}(\bs{\lm})}(X) = \nabla_{\bs{\theta}} \log p_{\bs{\theta}}(X)\Big|_{\bs{\theta}=\bs{\theta}(\bs{\lm})} \cdot D_{\bs{\lm}}\bs{\theta}(\bs{\lm}) = \dot{\ell}_{\bs{\theta}(\bs{\lm})}(X) \cdot D_{\bs{\lm}}\bs{\theta}(\bs{\lm}).
    \]
    Hence, the Fisher information transforms as a quadratic form:
    \[
    \bs{I}(\bs{\bs{\lm}}) \equiv \mathsf{Var}_{\bs{\bs{\lm}}}[\dot{\ell}_{\bs{\bs{\lm}}}(X)] = D_{\bs{\bs{\lm}}} \bs{\theta}(\bs{\bs{\lm}})^{\top} \bs{I}(\bs{\theta}(\bs{\bs{\lm}})) D_{\bs{\bs{\lm}}} \bs{\theta}(\bs{\bs{\lm}}).
    \]
\end{proof}

\begin{remark}
    The transformation of the Fisher information follows the logic of the chain rule. If one researcher parametrizes in terms of \(\bs{\theta}\) and another in terms of \(\bs{\lm}\), the information matrices are related simply by picking up the Jacobian of the change of parametrization. This changes the numerical values of the matrix entries (e.g., information about the standard deviation looks different than information about the variance), but it respects the geometry of the problem.
\end{remark}

\begin{ex}[Gamma Distribution Ambiguity]
    Why does this matter? Consider the Gamma distribution. Some software packages parametrize it using a \textit{rate} parameter \(\beta\) (density \(\propto e^{-\beta x}\)), while others use a \textit{scale} parameter \(1/\beta\) (density \(\propto e^{-x/\beta}\)). This is a diffeomorphic change of variables. While the definitions differ, the chain rule ensures we can always convert the Fisher information from one coordinate system to the other.
\end{ex}

Despite the information matrix changing forms, the fundamental limit on estimation accuracy remains constant.

\begin{pro}
    The Cram\'er--Rao bound is invariant under reparametrization.
\end{pro}

\begin{proof}
    In the original model \(\cc{P} = \{ \mathsf{P}_{\bs{\theta}} : \bs{\theta} \in \bs{\theta} \}\), the CR bound for estimating \(\bs{\theta}\) at the distribution given by \(\bs{\theta} = \bs{\theta}_0\) is \(\bs{I}(\bs{\theta}_0)^{-1}\).

    In the reparametrized model \(\cc{P}=\{\mathsf{P}_{\bs{\theta}(\bs{\lm})}:\bs{\lm}\in\bs{\lm}\}\), the considered distribution is \(\mathsf{P}_{\bs{\theta}_0}=\mathsf{P}_{\bs{\theta}(\bs{\lm}_0)}\) for \(\bs{\lm}_0:=\bs{\theta}^{-1}(\bs{\theta}_0)\). In the new parametrization, we estimate the function \(\gamma(\bs{\lm}) = \bs{\theta}(\bs{\lm})\). 
    
    By \autoref{pro:crb}, the CR bound for estimating this function is:
    \[
    \begin{split} 
    & D_{\bs{\lm}}\bs{\theta}(\bs{\lm}) \cdot \bs{I}(\bs{\lm})^{-1} \cdot D_{\bs{\lm}}\bs{\theta}(\bs{\lm})^{\top} \Big|_{\bs{\lm}=\bs{\lm}_{0}} \\ 
    & = D_{\bs{\lm}}\bs{\theta}(\bs{\lm}) \Big|_{\bs{\lm}=\bs{\lm}_{0}} \cdot \left( D_{\bs{\lm}}\bs{\theta}(\bs{\lm})^{\top} \Big|_{\bs{\lm}=\bs{\lm}_{0}} \cdot \bs{I}(\bs{\theta}_{0}) \cdot D_{\bs{\lm}}\bs{\theta}(\bs{\lm}) \Big|_{\bs{\lm}=\bs{\lm}_{0}} \right)^{-1} D_{\bs{\lm}}\bs{\theta}(\bs{\lm})^{\top} \Big|_{\bs{\lm}=\bs{\lm}_{0}} \\ 
    & = D_{\bs{\lm}}\bs{\theta}(\bs{\lm}) \cdot \left( D_{\bs{\lm}}\bs{\theta}(\bs{\lm})^{-1} \cdot \bs{I}(\bs{\theta}_{0})^{-1} \cdot (D_{\bs{\lm}}\bs{\theta}(\bs{\lm})^{\top})^{-1} \right) \cdot D_{\bs{\lm}}\bs{\theta}(\bs{\lm})^{\top} \\
    & = \bs{I}(\bs{\theta}_{0})^{-1}. 
    \end{split}
    \]
\end{proof}

\begin{note}
    Intuitively, this invariance \textit{must} hold. We are estimating the same characteristic of the same data-generating distribution \(\mathsf{P}_{\bs{\theta}_0}\). Whether we label the distribution using \(\bs{\theta}\) or \(\bs{\lm}\) is merely a naming convention. 
    
    For example, if the true data follows a Gaussian distribution with mean 5 and variance 7, the minimum variance for estimating the mean is a fixed number. Renaming the parameters does not change the difficulty of the statistical problem. The chain rule terms in the derivative of the target parameter \(\gamma(\bs{\lm})\) exactly cancel out the chain rule terms from the Fisher information transformation.
\end{note}
